\section{Fragments: SE₂, REC, RE, PP₂, SE₄}\label{sec:1}
\textbf{Setting.} Parameters and unknowns are integers. For $k, r ∈ ℕ$ write $\mathrm{Poly} k r := \mathrm{MvPolynomial} (\mathrm{Fin} k ⊕ \mathrm{Fin} r) ℤ$ and, for $P ∈ \mathrm{Poly} k r, a ∈ ℤ^{k}, u ∈ ℤ^{r}, \mathrm{evalAtZZ} P a u := P(a, u)$. A polynomial \emph{represents} $A ⊆ ℤ^{k}$ when $a ∈ A ⟺ ∃ u ∈ ℤ^{r}, P(a,u) = 0$ (\texttt{RepresentsZZ}).

\begin{definition}[\texttt{OneSortedFragments.\allowbreak{}lean}]\label[definition]{definition:1.1}
For $d, k ∈ ℕ$:

\begin{itemize}
\item $\mathrm{WithinSingleEquationZZ}\ d\ A :⟺ ∃ r, ∃ P ∈ \mathrm{Poly}\ k\ r, \mathrm{totalDegree}\ P ≤ d ∧ P \text{ represents } A$; $SE_{d}(k) := \{A \mid \mathrm{WithinSingleEquationZZ}\ d\ A\}$.
\item $\mathrm{WithinPPZZ}\ d\ A :⟺ ∃ r\ m, ∃ S : \mathrm{Fin}\ m → \mathrm{Poly}\ k\ r, (∀ i, \mathrm{totalDegree}\ (S\ i) ≤ d) ∧ (a ∈ A ⟺ ∃ u, ∀ i, S\ i\ (a,u) = 0)$; $PP_{d}(k) := \{A \mid \mathrm{WithinPPZZ}\ d\ A\}$.
\item $\mathrm{REC}(k) := \{A \mid \mathrm{ComputablePred}\ (· ∈ A)\}$, $\mathrm{RE}(k) := \{A \mid \mathrm{REPred}\ (· ∈ A)\}$ (mathlib \texttt{Mathlib.\allowbreak{}Computability}), through the \texttt{Primcodable (Fin k → ℤ)} instance.
\item $\mathrm{Dioph}(k) := \{A \mid ∃ r\ P, P \text{ represents } A\}$ (no degree bound).
\end{itemize}
\end{definition}
The number of unknowns r is existentially quantified and never bounded: the fragment is cut out by the \emph{degree} alone. This is the opposite convention to the existential-rank literature, which bounds the number of unknowns and leaves the degree free; the two conventions are orthogonal in intent. We prove no comparison between the two hierarchies in either direction, and none is claimed.

\begin{proposition}[unconditional lattice; \texttt{FragmentLattice.\allowbreak{}lean}, \texttt{singleEqFragmentZZ\_\allowbreak{}subset\_\allowbreak{}ppFragmentZZ}, \texttt{singleEqFragmentZZ\_\allowbreak{}subset\_\allowbreak{}diophFragmentZZ}, \texttt{fragment\_\allowbreak{}lattice\_\allowbreak{}ZZ}]\label[proposition]{proposition:1.2}
For every $k$ and every $d$: $SE_{d}(k) ⊆ PP_{d}(k) ⊆ \mathrm{Dioph}(k)$ (the two generic inclusions; \texttt{fragment\_\allowbreak{}lattice\_\allowbreak{}ZZ} is the $d = 2$ chain); $SE_{4}(k) = PP_{2}(k) = \mathrm{Dioph}(k)$ (\texttt{singleEqFragmentZZ\_\allowbreak{}four\_\allowbreak{}eq\_\allowbreak{}ppFragmentZZ\_\allowbreak{}two}, \texttt{singleEqFragmentZZ\_\allowbreak{}four\_\allowbreak{}eq\_\allowbreak{}diophFragmentZZ}).

\end{proposition}
\begin{proof}
A finite system $S_{1} = … = S_{m} = 0$ over ℤ has the same solutions as the single equation $ΣS_{i}^{2} = 0$, of degree at most 2d; with $d = 2$ this lands in SE₄. Conversely, Skolem's reduction replaces every monomial of degree ≥ 3 by a fresh unknown together with a quadratic defining equation, so any Diophantine set is PP₂; iterating the reduction terminates because each step strictly decreases the multiset of monomial degrees.

\end{proof}
\begin{proposition}[\texttt{ComputabilityFragments.\allowbreak{}lean}]\label[proposition]{proposition:1.3}
$SE_{d}(k) ⊆ \mathrm{RE}(k)$ and $PP_{d}(k) ⊆ \mathrm{RE}(k)$ for every $d$, unconditionally (\texttt{singleEquation\_\allowbreak{}relation\_\allowbreak{}re}, \texttt{ppRelation\_\allowbreak{}re}; the degree-four and degree-two cases are \texttt{singleQuartic\_\allowbreak{}relation\_\allowbreak{}re} and \texttt{ppQuadratic\_\allowbreak{}relation\_\allowbreak{}re}): witness search by \texttt{Nat.\allowbreak{}rfind} over a primitive-recursive solution checker (\texttt{primrec₂\_\allowbreak{}solCheck}). $\mathrm{REC}(k) ⊆ \mathrm{RE}(k)$ (\texttt{recFragmentZZ\_\allowbreak{}subset\_\allowbreak{}reFragmentZZ}).

\end{proposition}
\begin{proof}
For $SE_{d}$, let A be represented by $P ∈ \mathrm{Poly} k r$. The predicate $(a, u) ↦ (P(a,u) = 0)$ is primitive recursive: evaluating a fixed integer polynomial is a composition of additions and multiplications, and equality on ℤ is decidable. Since ℤʳ is \texttt{Primcodable}, the unbounded search \texttt{Nat.\allowbreak{}rfind} over the codes of witnesses halts on a exactly when some u is a root, so membership is r.e. For $PP_{d}$ the checker tests the finitely many equations of the system instead of one, which is again primitive recursive, and the same search applies. $\mathrm{REC}(k) ⊆ \mathrm{RE}(k)$ is the standard fact that a decidable predicate is r.e.
\end{proof}
\begin{proposition}[Arity zero; \texttt{ArityZeroFragments.\allowbreak{}lean}, \texttt{relation\_\allowbreak{}phase\_\allowbreak{}diagram\_\allowbreak{}arity\_\allowbreak{}zero}]\label[proposition]{proposition:1.4}
$ℤ^{0}$ is a singleton, so every relation is $∅$ or $\mathrm{univ}$; all five classes coincide and equal the full power set (\texttt{relation\_\allowbreak{}phase\_\allowbreak{}diagram\_\allowbreak{}arity\_\allowbreak{}zero}, unconditional). All strictness statements below are therefore stated for positive arity. This is not a triviality to be waved away: it is the reason every exactness statement in Sections 6 and 7 carries an arity hypothesis, and a phase diagram asserted "for all arities" would be false as stated.
\end{proposition}
\begin{proof}
ℤ⁰ has exactly one element, so $𝒫(ℤ^{0}) = \{∅, \mathrm{univ}\}$; univ is represented by the equation $0 = 0$ and ∅ by $1 = 0$, both of total degree at most two and with no witnesses, and both sets are decidable and r.e. So the five classes all equal the full power set, and no strictness statement can hold at arity zero.
\end{proof}

This paper has two load-bearing external inputs, together with one auxiliary proposition derived from them. Skolem's degree reduction is \emph{not} an external input: it is proved here (\cref{theorem:6.2}). The named propositions are the following. Each is a \texttt{def … : Prop} in the formalization and appears as an explicit hypothesis wherever it is used; none is an axiom, so no result silently depends on any of them.

\begin{definition}[Grunewald–Segal, code-uniform form; \texttt{QuadraticCodes.\allowbreak{}lean}]\label[definition]{definition:1.5}
A \emph{quadratic code} is a pair $c = (r, T)$ with $r ∈ ℕ$ the witness arity and T a finite list of terms (coefficient, at most two variable codes); it is an explicit finite datum, and the type of codes is \texttt{Primcodable}. The evaluation of a code at parameters a and an unknown list u reads variable code i as the i-th entry of the environment list ofFn a ++ u, out-of-range codes as 1, absent factors as 1 (\texttt{QuadraticCode.\allowbreak{}evalZZ}). The code c is \emph{solvable at a} if some u of length r is a root (\texttt{QuadraticCode.\allowbreak{}solvableZZ}). We set

\[
\mathrm{UniformQuadraticSolvabilityOnCodes} :⟺ ∀ k, \mathrm{ComputablePred}\ (\mathrm{fun}\ (c, a) ⇒ c \text{ solvable at } a).
\]
\end{definition}

The Grunewald–Segal procedure yields this hypothesis after an effective translation from our finite quadratic code to the coefficient-table representation it takes as input; what \cref{definition:1.5} records is that coding-equivalent consequence, not their statement. The shape is the same on both sides — a coefficient table and a witness arity in, a yes/no answer out, uniformly in the table — but we do not assert that the two inputs are the same object. The citation needs care. The 1981 note \cite{GS81} is, in its authors' own description, a sketch of an effective procedure for an arbitrary quadratic equation with rational coefficients; the algorithm is presented in full in \cite{GS04}, and it is there that the input is stated to be an equation \emph{in arbitrarily many unknowns}, which is the generality our fragment requires. We therefore cite \cite{GS04} as the source of \cref{definition:1.5} and \cite{GS81} as its announcement. Two things should be said precisely about this hypothesis. First, we do not claim that our coding is theirs: \cref{theorem:6.4} supplies the compiler that turns an instance in our form into a parameter-free coefficient table, and \cref{theorem:6.4}(f) shows that this compiler is primitive recursive, so the \emph{translation} of inputs is effective. Second, the hypothesis is uniformity of the decision procedure in the table. We isolate this code-uniform quadratic-solvability principle as an explicit external hypothesis. The published abstract of \cite{GS04} describes an algorithm that takes an arbitrary quadratic equation as its input, which supports the required generality; we have not audited the body of \cite{GS04} against \cref{definition:1.5}, and every result below that depends on the definition remains explicitly conditional on the principle as stated there. \cref{theorem:6.4}(f) does not establish uniformity and is not asked to; it establishes that the translation between our instance form and that input is effective. We therefore state \cref{definition:1.5} as the classical input in its uniform form, and we do not claim to have derived uniformity from per-equation decidability.

Three conventions must be matched before the two statements can be compared at all, and we record how each is handled. \emph{Coefficients.} \cite{GS04} admits rational coefficients; a quadratic code carries integer ones. Multiplying an equation through by a common denominator changes neither its solution set nor its total degree, so the instances we compile are a subclass of theirs and no clearing step is needed in the direction we use. \emph{Unknowns.} Ours range over ℤ. The headline statement of \cite{GS04} concerns solutions in positive integers, so the convention must be addressed rather than assumed away, and it is addressed by \cite{GS04} itself: the algorithm there is stated for integral solutions constrained to a region bounded by an arbitrary finite collection of rational affine hyperplanes, and the empty collection gives unrestricted integral solutions. The ℤ form is thus part of what is claimed, not something we reconstruct. \cite{GS81} states the integer case directly, but describes itself as a sketch, which is why we cite \cite{GS04} as the source. Independently of both, the substitution $u = u^{+} − u^{-}$ with $u^{+}, u^{-} ≥ 1$ carries integer solutions to positive ones and back, doubling the number of unknowns and leaving the degree at two; the number of unknowns is unrestricted in our fragment, so nothing we assert would fail under the other convention either. \emph{Arity.} The witness arity is part of the code, so within each fixed parameter arity k the decider of \cref{definition:1.5} is uniform in the code and hence in the witness arity. The ∀ k of \cref{definition:1.5} quantifies over the parameter arity outside \textrm{ComputablePred}: it asserts, for each fixed k, one decider uniform in the finite quadratic code and the parameter tuple, and it does not assert a single algorithm uniform in k. The development also formalizes the global form — one decider on the sigma-coded instance type of codes paired with parameter lists (\texttt{GlobalQuadraticSolvabilityDecidable}) — and proves that it implies \cref{definition:1.5} (\texttt{globalGS\_\allowbreak{}implies\_\allowbreak{}uniformGS}); the hypothesis used in this paper is therefore the weaker, fixed-arity form, and every result conditional on it holds a fortiori under the global principle.

The clause-by-clause crosswalk from \cite{GS04} to \cref{definition:1.5} is \texttt{literature/GS\_\allowbreak{}CROSSWALK.\allowbreak{}md} in the artifact, with the level of evidence recorded row by row. One point from it belongs here. \cite{GS04} states a decision procedure, not a bound: it lists a finite solution set and reports infinitude otherwise, and its auxiliary result on the density at infinity of the translates of a rational point is qualitative. That reading is corroborated from outside: \cite{Die03}, appearing a year earlier, records that for $s = 3$ and $s = 4$ no polynomial bounds were then known. So the non-effectivity in \cref{remark:2.3} is not a bound we have declined to look up; it reflects the state of the line we are citing.

\begin{definition}[DPRM over ℕ; \texttt{DPRMTransport.\allowbreak{}lean}]\label[definition]{definition:1.6}
$\mathrm{WithinDiophNN}\ A :⟺ ∃ s, ∃ P ∈ \mathrm{Poly}\ k\ s, ∀ a ∈ ℕ^{k}, a ∈ A ⟺ ∃ y ∈ ℕ^{s}, P(a,y) = 0$ (integer coefficients, natural unknowns); $DPRM_{NN} :⟺ ∀ k\ A, \mathrm{REPred}\ (· ∈ A) → \mathrm{WithinDiophNN}\ A$. This is the theorem of \cite{DPR61,Mat70} in its standard form.

\end{definition}
\begin{definition}[auxiliary, derived]\label[definition]{definition:1.7}
QuadraticSolvabilityDecidableZZ (per-equation decidability) and $DPRM_{ZZ}$ (every r.e. relation on ℤᵏ is Dioph) are \emph{consequences} of \cref{definition:1.5,definition:1.6} (\cref{theorem:6.4,theorem:6.7}); they are kept as intermediate hypotheses so that each result can name the weakest input it actually uses.

\end{definition}
